\section{Conclusiones}
\label{sec:conclusiones}

El presente trabajo abordó el problema de la toma de decisiones en el juego Tetris mediante técnicas de Aprendizaje por Refuerzo y optimización estocástica. El entorno fue modelado formalmente como un Proceso de Decisión de Markov, permitiendo representar el problema dentro del marco teórico del Aprendizaje por Refuerzo.

Con el fin de reducir la complejidad del espacio de estados, se empleó una representación basada en características geométricas del tablero, entre las que destacan la altura agregada, el número de huecos, la rugosidad de la superficie, las líneas eliminadas, la altura máxima y los pozos.

La función de valor fue aproximada mediante un modelo lineal cuyos parámetros se optimizaron utilizando dos estrategias diferentes: Aprendizaje por Diferencias Temporales TD(0) y el Método de Entropía Cruzada.

Los experimentos realizados mediante un barrido exhaustivo de hiperparámetros permitieron identificar una configuración de entrenamiento capaz de eliminar en promedio

\begin{equation}
\bar{L}=409.13
\end{equation}

líneas antes de alcanzar el estado terminal, evidenciando la capacidad del agente para aprender políticas de juego competitivas.

Más allá del barrido, la evaluación controlada entre ambos métodos reveló que bajo horizontes de juego cortos, las políticas exhiben rendimientos similares debido a un efecto techo. Sin embargo, al extender la evaluación a $5000$ colocaciones para medir la verdadera capacidad de supervivencia, el Método de Entropía Cruzada supera ampliamente a TD(0) ($1048.2$ frente a $563.2$ líneas promedio). Esta diferencia sustancial se explica por la estabilidad paramétrica: mientras los pesos de TD(0) divergen significativamente manifestando la inestabilidad propia del aprendizaje semi-gradiente e induciendo un estado terminal prematuro, CEM mantiene un vector de pesos acotado y estrictamente normalizado ($\lVert w \rVert = 1$). Esto garantiza un control robusto a largo plazo, consolidando a la optimización estocástica directa como una alternativa netamente superior al bootstrapping en este contexto.

Los resultados obtenidos muestran que las técnicas de aproximación lineal de funciones constituyen una alternativa efectiva para abordar problemas secuenciales de gran complejidad sin necesidad de recurrir a modelos de Aprendizaje Profundo o arquitecturas neuronales de gran escala.

Finalmente, la naturaleza modular de la arquitectura desarrollada permite la incorporación de nuevos algoritmos de aprendizaje, convirtiendo al sistema propuesto en una plataforma adecuada para futuras investigaciones en Aprendizaje por Refuerzo aplicado a videojuegos y problemas de optimización secuencial.